\chapter{Практика}

\section{Настройка имени}

\begin{lstlisting}[language=C, caption={Пример настройки имени устройства}]
# переход в режим конфигурации
configure
# задаем имя
hostname R1
# принимаем настройки
do commit
do confirm
\end{lstlisting}

\section{Настройка времени}

\begin{lstlisting}[language=C, caption={Пример настройки времени}]
# В режиме пользователя
set date 16:49:00 25 March 2026
# вносим изменения
commit
# проверка
show date
# переходим в режим конфигурации
configure
# задаем часовой пояс
clock timezone gmt +7
# вносим изменения
do commit
do confirm
\end{lstlisting}

\section{Настройка IP на интерфейсе}

\begin{lstlisting}[language=C, caption={Пример настройки IP на интерфейсе}]
# переход в режим конфигурации
configure
# переход к настройке интерфейса gi 1/0/1
interface gigabitethernet 1/0/1
# задаем ip с маской
ip address 192. 168. 100. 1/24
# отключаем firewall
ip firewall disable
# подтверждаем изменения
do commit
do confirm
\end{lstlisting}

\section{Настройка IP на интерфейсе}

\begin{lstlisting}[language=C, caption={Пример настройки IP на интерфейсе}]
# переход в режим конфигурации
configure
# переход к настройке интерфейса gi 1/0/1
interface gigabitethernet 1/0/1
# задаем ip с маской
ip address 192. 168. 100. 1/24
# отключаем firewall
ip firewall disable
# подтверждаем изменения
do commit
do confirm
\end{lstlisting}

\section{Создание и настройка Loopback интерфейса}

\begin{lstlisting}[language=C, caption={Пример создания и настройки Loopback интерфейса}]
# переход в режим конфигурации
configure
# переход к настройке интерфейса Loopback 1
interface Loopback 1
# задаем ip с маской
ip address 8.8.8.8/32
exit
do commit
do confirm
\end{lstlisting}

\section{Создание и настройка sub интерфейса}

\begin{lstlisting}[language=C, caption={Пример создания и настройки sub интерфейса}]
configure
# переходим в настройку интерфейса gi 1/0/1. Через точку указан vlan для sub интерфейса
interface gigabitethernet 1/0/1.50
# задаем ip
ip address 172. 16. 25. 1/24
# отключаем firewall
ip firewall disable
\end{lstlisting}

\section{Создание и настройка пользователя}

\begin{lstlisting}[language=C, caption={Пример создания и настройки пользователя}]
configure
# создаем пользователя master и переходим к его настройке
username master
# задаем пароль
password eltex
# задаем уровень привилегий
privilege 14
\end{lstlisting}

\section{Включение ssh сервера}

\begin{lstlisting}[language=C, caption={Пример включения ssh сервера}]
configure
ip ssh server
\end{lstlisting}


\section{Настройка IP на ПК (Debian)}

Чтобы посомтреть доступные интерфейсы и их параметры необходимо выполнить команду ip a:

\begin{figure}[H] 
    \centering
    \includegraphics[width=0.75\textwidth]{linux.png}
    \caption{Результат работы команды}
\end{figure}

Настройка интерфейсов проиводится в файле /etc/network/interfaces. Файл необходимо отредактировать следующим образом:

\begin{figure}[H] 
    \centering
    \includegraphics[width=0.75\textwidth]{linux.png}
    \caption{Настройка интерфейсов}
\end{figure}

Проверим наличие связи между ПК и роутером:

\begin{figure}[H] 
    \centering
    \includegraphics[width=0.75\textwidth]{linux2.png}
    \caption{Проверка сетевой связности}
\end{figure}

БАГ: почему-то связь устанавливается только при подключении к интерфейсу gi 1/0/3, хотя на роутере настраивали gi 1/0/1.

\section{Подключение к роутеру по ssh}

Для подключения по ssh необходимо в терминале ввести команду ssh <username>@<ip\_address>. В нашем случае необходимо
ввести команду ssh admin@192.168.100.1

\begin{figure}[H] 
    \centering
    
\includegraphics[width=0.75\textwidth]{linux3.png}
    \caption{Подключение по ssh}
\end{figure}

\section{Настройка баннеров}

\begin{lstlisting}[language=C, caption={Пример включения ssh сервера}]
configure
# баннер перед входом
banner login Router-R1
# баннер после входа
banner exec «Connect Router R1»
do commit
do confirm
\end{lstlisting}

\begin{figure}[H] 
    \centering
    
\includegraphics[width=0.75\textwidth]{banners.png}
    \caption{Пример работы баннеров}
\end{figure}

\endinput